\section{Задание}
\begin{bbttask}
  \item собрать в симуляторе топологию PC1 — R1 — R2 — PC2, привести схему с именами
  интерфейсов;\label{task:topology}
  \item назначить адреса интерфейсам маршрутизаторов и компьютеров по плану адресации своего
  варианта, свести адреса в таблицу;\label{task:address}
  \item проверить связь PC1 с PC2 до настройки маршрутов и объяснить
  результат;\label{task:before}
  \item настроить на R1 и R2 статические маршруты в чужую локальную сеть, привести таблицу
  маршрутизации R1;\label{task:routes}
  \item проверить связь командами \texttt{ping} и \texttt{tracert} и объяснить путь
  пакетов.\label{task:check}
\end{bbttask}

План адресации варианта~\bbtVariant{} приведён в таблице~\ref{tab:plan}.

\begin{bbttable}{План адресации варианта}{tab:plan}
  \begin{tabular}{|l|l|}
    \hline
    Сеть & Адрес \\ \hline
    LAN1 (PC1 — R1) & 192.168.5.0/24 \\ \hline
    LAN2 (R2 — PC2) & 192.168.15.0/24 \\ \hline
    Между R1 и R2 & 10.0.5.0/30 \\ \hline
  \end{tabular}
\end{bbttable}
